\begin{tikzpicture}[
  >={Latex[length=3mm,width=2.2mm]},
  every node/.style={font=\small},
  uvec/.style={->,line width=1pt,black},
  ivec/.style={->,line width=1.2pt,vpurple},
  bvec/.style={->,line width=1.3pt,bblue},
  dashline/.style={dashed,thin,ered}
]

% Piano in prospettiva (giallo chiaro)
\fill[ytangent,opacity=0.4,draw=ytangent!70!black]
  (-3,-1.5) -- (4,-1.0) -- (4,1.2) -- (-3,0.7) -- cycle;

% Filo rettilineo verticale (grigio)
\fill[blobgray!60,draw=black!50] (-0.15,-2.5) rectangle (0.15,3.0);

% Corrente i (viola, dal basso)
\draw[ivec] (-0.3,-2.0) -- ++(0,0.8) node[left] {$i$};

% ds in alto
\draw[uvec,line width=1.2pt] (0,2.0) -- ++(0,1.0) node[above] {$d\mathbf{s}$};

% Angolo theta
\draw[thin] (0,2.0) ++(0.5,0) arc (0:35:0.5);
\node at (0.65,2.25) {$\theta$};

% Raggio 2a (cerchio tratteggiato rosso nel piano)
\draw[dashline] (0,0) ellipse (2.2 and 0.5);
\node[left] at (-2.2,0) {$2a$};

% R orizzontale
\draw[thin] (0,0.1) -- (2.2,0.1);
\node[above] at (1.1,0.1) {$R$};

% Punto P
\coordinate (P) at (2.5,-0.1);
\fill (P) circle (2pt);
\node[below right] at (P) {$P$};

% u_r da ds verso P
\draw[uvec] (0.3,2.0) -- ++(1.2,-0.8) node[below] {$\mathbf{u}_r$};

% r da ds a P
\draw[thin] (0,2.0) -- (P);
\node[right] at (1.2,1.0) {$r$};

% dB (azzurro, verso destra in basso)
\draw[bvec] (P) -- ++(1.0,0.4) node[right] {$d\mathbf{B}$};

\end{tikzpicture}
